% Part: lambda-calculus
% Chapter: lambda-definability
% Section: truth-values

\documentclass[../../../include/open-logic-section]{subfiles}

\begin{document}

\olfileid{lam}{ldf}{tvr}
\olsection{సత్యమూల్యాలు మరియు సంబంధాలు}

శుద్ధ లాంబ్డా కలనశాస్త్రంలో సత్యమూల్యాలను ఇలా
సంకేతీకరించవచ్చు:
\begin{align*}
  \fn{true} & \ident \lambd[x][\lambd[y][x]]\\
  \fn{false} & \ident \lambd[x][\lambd[y][y]]
\end{align*}

సత్యమూల్యాలను \emph{ఎంపిక ప్రమేయాలు}గా సూచిస్తున్నాం:
అవి రెండు ఆర్గ్యుమెంట్లను స్వీకరించి, వాటిలో ఒకదాన్ని
తిరిగి ఇచ్చే ప్రమేయాలు. $\fn{true}$ మొదటి ఆర్గ్యుమెంటును,
$\fn{false}$ రెండవదాన్ని ఎంచుకుంటుంది. ఉదాహరణకు,
$\fn{true}\, M N$ ఎల్లప్పుడూ $M$కు తగ్గుతుంది;
$\fn{false}\, M N$ ఎల్లప్పుడూ $N$కు తగ్గుతుంది.

\begin{defn}
$R \subseteq \Nat^k$ అనే సంబంధానికి, ఒక పదం~$R$ ఉండి,
\begin{align*}
  R\, \num{n_1} \dots \num{n_k} & \bred \fn{true}
  \intertext{$R(n_1, \dots, n_k)$ అయినప్పుడల్లా; అలాగే సంబంధం వర్తించని ఇతర సందర్భాల్లో}
  R\, \num{n_1} \dots \num{n_k} & \bred \fn{false}
\end{align*}
అయితే, ఆ సంబంధాన్ని
\tetoken{లాంబ్డాతో నిర్వచించదగిన}{!!{lambda definable}} సంబంధం
అంటాం.
\end{defn}
\sourcecorrection{OLTELAMLDFTVR-001}{మూలం సంబంధాన్ని $R\subseteq\Nat^n$గా ప్రకటించి, తదుపరి రెండు సూత్రాల్లో $n_1,\dots,n_k$ అనే $k$ ఆర్గ్యుమెంట్లను వాడింది. ఆ సూత్రాలు, వాటి షరతును యథాతథంగా ఉంచి, సంబంధపు స్థానం $k$ అని చూపడానికి ప్రకటించిన ఘాతాన్ని $\Nat^k$గా సరిచేశాం.}

ఉదాహరణకు, $0$కు వర్తించి, $0$కు మాత్రమే వర్తించే
$\fn{IsZero} = \{0\}$ సంబంధం ఈ పదం ద్వారా
\tetoken{లాంబ్డాతో నిర్వచించదగినది}{!!{lambda definable}}:
\[
  \fn{IsZero} \ident \lambd[n][n (\lambd[x][\fn{false}])\, \fn{true}].
\]
ఇది ఎలా పనిచేస్తుంది? చర్చ్ సంఖ్యాంకాలు పునరావర్తకాలుగా
నిర్వచించబడ్డాయి: అవి మొదటి ఆర్గ్యుమెంటును రెండవదానికి
$n$ సార్లు ప్రయోగిస్తాయి. మొదటి విలువను $\fn{true}$గా
పెడతాం. పునరావర్తనంలోని ప్రతి దశలో, గత దశ ఫలితం
ఏదైనా సరే $\fn{false}$ను తిరిగి ఇస్తాం. ఈ దశ మొదటి
విలువకు $n$ సార్లు వర్తిస్తుంది. అది అసలు వర్తించకపోతేనే,
అంటే $n=0$ అయితేనే, ఫలితం $\fn{true}$ అవుతుంది.

ఈ సత్యమూల్య ప్రాతినిధ్యం ఆధారంగా కొన్ని సత్య ప్రమేయాలను
కూడా నిర్వచించవచ్చు. నిషేధం, సంయోగం ప్రాతినిధ్యాలు ఇవి:
\begin{align*}
  \fn{Not} & \ident \lambd[x][x\, \fn{false}\, \fn{true}]\\
  \fn{And} & \ident \lambd[x][\lambd[y][xy \,\fn{false}]]
\end{align*}
$\fn{Not}$ ఒక ఆర్గ్యుమెంటును స్వీకరించి, అది
$\fn{false}$ అయితే $\fn{true}$ను, అది $\fn{true}$
అయితే $\fn{false}$ను తిరిగి ఇస్తుంది. $\fn{And}$ రెండు
సత్యమూల్యాలను ఆర్గ్యుమెంట్లుగా స్వీకరిస్తుంది; రెండూ
$\fn{true}$ అయినప్పుడే $\fn{true}$ను తిరిగి ఇవ్వాలి.
పైన వివరించినట్లు సత్యమూల్యాలు ఎంపిక ప్రమేయాలుగా
సూచించబడ్డాయి. కాబట్టి $x$ సత్యమూల్యమై, దాన్ని రెండు
ఆర్గ్యుమెంట్లకు ప్రయోగిస్తే, $x$ అనేది $\fn{true}$ అయితే మొదటి
ఆర్గ్యుమెంటు, లేకపోతే రెండవది ఫలితం. ఇప్పుడు
$\fn{And}$ తన రెండు ఆర్గ్యుమెంట్లు $x$, $y$లను
స్వీకరించి, మొదటి ఆర్గ్యుమెంటు~$x$కు $y$,
$\fn{false}$లను అందిస్తుంది. $x$ సత్యమూల్యమని
అనుకుంటే, $x$ అనేది $\fn{true}$ అయినప్పుడు ఫలితం~$y$;
$x$ అనేది $\fn{false}$ అయినప్పుడు ఫలితం~$\fn{false}$.
కావలసింది ఇదే.

ఇక్కడ $\fn{And}$కు ఆర్గ్యుమెంట్లుగా సత్యమూల్యాలను మాత్రమే
ఇస్తున్నామని గమనించండి. ఇతర పదాలను ఇస్తే, ఫలితం
(నియత రూపం ఉంటే అది) తప్పనిసరిగా సత్యమూల్యం కాకపోవచ్చు.

\begin{prob}
  సత్యమూల్యాలను $\lambd$-పదాలుగా సంకేతీకరించే ఈ పద్ధతిని
  ఉపయోగించి, కనీసం ఒకటి నిజమైతే నిజమయ్యే సమావేషక
  వికల్పం, కచ్చితంగా ఒకటి మాత్రమే నిజమైతే నిజమయ్యే
  బహిష్కార వికల్పం అనే సత్య ప్రమేయాలను సూచించే
  $\fn{Or}$, $\fn{Xor}$
  ప్రమేయాలను నిర్వచించండి.
\end{prob}

\end{document}
